\section{Proofs}\label{app:proofs}
% =============================================================================================
% Owner: proofs writer. Proofs of the statements of Section 4 (sections/theory.tex) and of
% Section 8.1 (sec:mk-theory), in the order of the main text. Source: docs/theory.md; the proofs
% of prop:anchored and prop:anchored-target are new. Checked against src/synthval
% (critic.time_block_labels, train_eval_split, perm_mean_diff, perm_pvalues, hac_var_mean,
% relative.relative_test_anchored).
% New labels: lem:exchangeable, lem:rank, lem:uniform-lip (registry), lem:kantorovich-ineq.
% Numerical constants of the examples were recomputed with a script (see the writer's report).
% =============================================================================================

This appendix proves the results of \cref{sec:theory} and \cref{sec:mk-theory} in the order of
the main text; \cref{prop:bayes-logit}, \cref{cor:jeffreys}, \cref{cor:block-bias},
\cref{prop:verdict} and \cref{rem:combine-seeds} are proved where they are stated. Statuses are those of
\cref{tab:theory-status}. Where a statement concerns the implemented procedure, the proof says
which features of the code it uses and what changes when the code departs from the assumptions.
Conditioning on $\Gsig$ fixes the critic, which is jointly measurable in the data and its
parameters.

% ---------------------------------------------------------------------------------------------
\subsection{What a window test can detect}\label{app:proofs-windows}

\begin{proof}[Proof of \cref{prop:window-ident}]
(a) Let $t_1<\dots<t_k$ with $t_k-t_1\le L-1$. By strict stationarity,
$(r_{t_1},\dots,r_{t_k})\eqd(r_1,r_{1+t_2-t_1},\dots,r_{1+t_k-t_1})$, and since $1+t_k-t_1\le L$
the right-hand side is a coordinate projection of $(r_1,\dots,r_L)\sim P_L$. Autocovariances,
moments $\E g(r_t,r_{t+h})$ with $h\le L-1$ and the law of $\sum_{t=1}^hr_t$ with $h\le L$ are
functionals of such joint laws.
(b) Every finite-dimensional distribution is, by (a), a marginal of $P_L$ for $L\ge t_k-t_1+1$.
Cylinder sets form a $\pi$-system that generates the product $\sigma$-field, so two process laws
with the same finite-dimensional distributions coincide (Dynkin's $\pi$-$\lambda$ theorem; this
is the uniqueness part of Kolmogorov's extension theorem).
(c) Let $(e_t)$ be \iid{} $N(0,1)$, $\theta\neq0$, and $r_t=(e_t+\theta e_{t-L})/\sqrt{1+\theta^2}$.
This stationary Gaussian process has $\gamma(0)=1$, $\gamma(h)=0$ for $1\le h\le L-1$ and
$\gamma(L)=\theta/(1+\theta^2)\neq0$. A Gaussian law is determined by its covariance, so
$P_L=N(0,I_L)$, as for \iid{} $N(0,1)$ returns, whereas the two laws $P_{L+1}$ differ in
$\Cov(r_1,r_{L+1})$.
\end{proof}

\begin{proof}[Proof of \cref{prop:markov-pairs}]
Let $\mu$ be the common law of $r_0$ and $r'_0$, and $K,K'$ the transition kernels, so that
$(r_0,r_1)\sim\mu\otimes K$ and $(r'_0,r'_1)\sim\mu\otimes K'$. Equality of the pair laws gives
$\int_AK(x,C)\,\mu(dx)=\int_AK'(x,C)\,\mu(dx)$ for all Borel $A,C$, hence $K(x,C)=K'(x,C)$ for
$\mu$-a.e.\ $x$ and each fixed $C$. The Borel $\sigma$-field of a Polish space is generated by a
countable $\pi$-system $\mathcal C$; the union over $C\in\mathcal C$ of the exceptional sets is a
$\mu$-null set $N$, and off $N$ the two probability measures $K(x,\cdot)$, $K'(x,\cdot)$ agree on
$\mathcal C$, hence everywhere. By the Markov property and stationarity,
$(r_0,\dots,r_k)\sim\mu(dx_0)K(x_0,dx_1)\cdots K(x_{k-1},dx_k)$. Each $x_j$ is integrated
against a measure with marginal $\mu$, which does not charge $N$, so replacing $K$ by $K'$ leaves
this law unchanged for every $k$. All finite-dimensional laws agree, and \cref{prop:window-ident}(b)
concludes. For a chain of order $k$, $s_t=(r_t,\dots,r_{t+k-1})$ is a stationary first-order
chain on the Polish space $S^k$ whose pair law is the law of $(r_0,\dots,r_k)$, that is
$P_{k+1}$; apply the first part to $(s_t)$.
\end{proof}

\begin{proof}[Sketch: \GARCHt{} returns are not Markov of any finite order]
Write $\varepsilon_t=r_t-\mu=\sigma_tz_t$ with $\sigma_{t+1}^2=\omega+a\varepsilon_t^2+b\sigma_t^2$,
$a,b>0$, $z_t$ \iid{} with an everywhere positive density, and let $\mathcal F_t=\sigma(r_s,s\le t)$.
In the stationary solution $\sigma_{t+1}^2$ is $\mathcal F_t$-measurable and
$\sigma_{t+1}^2=\E[\varepsilon_{t+1}^2\mid\mathcal F_t]$. Suppose the returns were Markov of order
$k$. Then the conditional law of $r_{t+1}$ given $\mathcal F_t$ depends only on
$W=(r_{t-k+1},\dots,r_t)$, so $\sigma_{t+1}^2=\phi_0(W)$ a.s. Iterating the recursion $k$ times,
\[
  b^k\sigma_{t-k+1}^2=\sigma_{t+1}^2-\omega\sum_{j<k}b^j-a\sum_{j<k}b^j\varepsilon_{t-j}^2 ,
\]
so $V:=\sigma_{t-k+1}^2=\phi(W)$ a.s.\ for a measurable $\phi$ (here $b>0$ is used). Now $V$ is
$\mathcal F_{t-k}$-measurable and $z_{t-k+1},\dots,z_t$ are independent of $\mathcal F_{t-k}$.
Given $V=v$, $W$ is the image of $(z_{t-k+1},\dots,z_t)$ under a triangular map whose Jacobian is
$\prod_s\sigma_s>0$, so the conditional law of $W$ given $V=v$ has an everywhere positive Lebesgue
density on $\R^k$ for every $v>0$. If $V=\phi(W)$, this conditional law is carried by the level
set $\{\phi=v\}$, which therefore has positive Lebesgue measure for almost every $v$ in the law
of $V$. Disjoint sets of positive Lebesgue measure are at most countably many, so the law of $V$
would be purely atomic. But given $\mathcal F_{t-k-1}$, $V=\omega+b\sigma_{t-k}^2+a\sigma_{t-k}^2z_{t-k}^2$
has a density (here $a>0$ is used), so the law of $V$ has no atoms: a contradiction. With $b=0$
(ARCH(1)) the returns are first-order Markov, so the condition $b>0$ cannot be dropped.
\end{proof}

% ---------------------------------------------------------------------------------------------
\subsection{The held-out critic statistic}\label{app:proofs-critic}

\begin{proof}[Proof of \cref{prop:unbiased}]
Conditionally on $\Gsig$, $\fhat=\ghat\circ A$ is a fixed Lipschitz function, and
$\abs{\fhat(x)}\le\abs{\fhat(0)}+\Lip(\fhat)\norm{x}$ makes it $P$- and $Q$-integrable for
$P,Q\in\Pspace{1}(\R^d)$.
(a) Each evaluation window has conditional law $P$ (resp.\ $Q$) given $\Gsig$, so
$\E[\fhat(X_i)\mid\Gsig]=P\fhat$ and $\E[\fhat(Y_j)\mid\Gsig]=Q\fhat$; average over $\Iev$ and
$\Jev$. The change of variables $u=A(x)$ gives $P\fhat=P_A\ghat$ and $Q\fhat=Q_A\ghat$.
(b) $P_A,Q_A\in\Pspace{1}$ because $A$ is affine, and $P_A(S_A)=Q_A(S_A)=1$ because supports have
full measure in $\R^d$. If $\ell:=\Lip_{S_A}(\ghat)=0$, $\ghat$ is constant on $S_A$ and
$\Dcrit(\fhat)=0$. Otherwise $h=\ghat/\ell$ is 1-Lipschitz on the closed set $S_A$, and McShane's
extension $\bar h(x)=\inf_{y\in S_A}\{h(y)+\norm{x-y}\}$ is finite (because
$h(y)+\norm{x-y}\ge h(y_0)-\norm{x-y_0}$ for a fixed $y_0\in S_A$), 1-Lipschitz on $\R^d$ and
equal to $h$ on $S_A$. Hence $\Dcrit(\fhat)=\ell\,(P_A\bar h-Q_A\bar h)\le\ell\,\W(P_A,Q_A)$ by
\cref{fact:kr-duality}, and $\ell\le\Lip(\ghat)$.
(c) Conditional independence of the evaluation windows gives the variance identity. For
independent $U,U'\sim P_A$ and Lipschitz $g$,
$\Var g(U)=\tfrac12\E\{g(U)-g(U')\}^2\le\tfrac12\Lip(g)^2\E\norm{U-U'}^2=\Lip(g)^2\Tr\Sigma_{P_A}$;
apply it with $g=\ghat$, using $\Var_P\fhat=\Var_{P_A}\ghat$ and
$\Sigma_{P_A}=\Lambda^{-1}\Sigma_P\Lambda^{-1}$, and likewise for $Q$.
\end{proof}

% ---------------------------------------------------------------------------------------------
\subsection{Exact level under exchangeability}\label{app:proofs-level}

We work conditionally on $\Gsig$, which fixes $\fhat$ and hence the map $t$ of
\cref{sec:theory-level}. Write $N_e=n_e+m_e$, $\mathcal K$ for the group of permutations of
$\{1,\dots,N_e\}$ and $\pi_0=\mathrm{id}$.

\begin{lemma}[Random permutations give an exchangeable vector]\label{lem:exchangeable}
Under (A1$'$) and (A1)(iv), the vector $(T_0,T_1,\dots,T_B)$ with $T_b=t(\pi_b\cdot V)$ is
exchangeable conditionally on $\Gsig$.
\end{lemma}

\begin{proof}
Let $\varpi$ be uniform on $\mathcal K$, independent of $(V,\pi_1,\dots,\pi_B)$. For every fixed
$g\in\mathcal K$, $g\cdot V\eqd V$ by (A1$'$), and $V$ is independent of $(\varpi,\pi_{1:B})$; hence
$(\varpi\cdot V,\pi_{1:B})\eqd(V,\pi_{1:B})$ and
\[
  \big(t(\pi_b\cdot V)\big)_{b=0}^B\eqd\big(t(\pi_b\cdot(\varpi\cdot V))\big)_{b=0}^B
  =\big(t((\pi_b\varpi)\cdot V)\big)_{b=0}^B ,
\]
because $(\pi,v)\mapsto\pi\cdot v$ is a group action. Given $\varpi$, the elements
$\pi_0\varpi,\dots,\pi_B\varpi$ are \iid{} uniform on $\mathcal K$ (right multiplication is a
bijection of $\mathcal K$ and $\pi_0\varpi=\varpi$ is uniform), so the tuple is \iid{} uniform and
independent of $V$. An \iid{} tuple is exchangeable, and applying the same map
$\sigma\mapsto t(\sigma\cdot V)$ to each coordinate preserves exchangeability.
\end{proof}

\begin{lemma}[Rank bound]\label{lem:rank}
If $(T_0,\dots,T_B)$ is exchangeable and $R=\#\{b\in\{0,\dots,B\}:T_b\ge T_0\}$, then
$\Prob(R\le k)\le k/(B+1)$ for $k=0,1,\dots,B+1$.
\end{lemma}

\begin{proof}
Let $R_c=\#\{b:T_b\ge T_c\}$. By exchangeability $R_c\eqd R_0=R$ for every $c$, so
$\Prob(R\le k)=(B+1)^{-1}\E\,\#\{c:R_c\le k\}$. Deterministically $\#\{c:R_c\le k\}\le k$: if the
set $C$ of such $c$ had more than $k$ elements, its element $c^*$ with the smallest $T_{c^*}$
(ties broken arbitrarily) would satisfy $R_{c^*}\ge\abs{C}>k$.
\end{proof}

\begin{proof}[Proof of \cref{prop:exact-level}]
$\pperm=R/(B+1)$ with $R$ as in \cref{lem:rank}, so
$\{\pperm\le\alpha\}=\{R\le\lfloor\alpha(B+1)\rfloor\}$. Apply \cref{lem:exchangeable,lem:rank}
conditionally on $\Gsig$ and take expectations. That (A1) and $H_0$ imply (A1$'$): by (A1)(ii)
the index sets are independent of the window values, so given $\Gsig$ the evaluation windows are
the windows not used for training, which are independent of the training windows and of $\xi$
by (A1)(i) and (iii); they are therefore conditionally \iid{} with law $P=Q$, and an \iid{}
vector is exchangeable.
\end{proof}

The code is covered in three respects. The default, conservative p-value (\texttt{perm\_pvalues})
counts the $b$ with $T_b\ge T-10^{-12}\max(1,\abs T)$, which can only enlarge $R$, so its size is
at most $\lfloor\alpha(B+1)\rfloor/(B+1)$, which is $50/1001\approx0.04995$ at $\alpha=0.05$ and
$B=1000$. The randomised p-value is $p=(G+U(1+E))/(1+B)$, where $G=\#\{b\ge1:T_b>T_0\}$,
$E=\#\{b\ge1:T_b=T_0\}$ and $U\sim U(0,1)$ is independent of everything else; it is exactly
$U(0,1)$ under (A1$'$) and (A1)(iv). Given $(T_0,\dots,T_B)$, the integer part $J$ of $U(1+E)$ is
uniform on $\{0,\dots,E\}$ and its fractional part $U'$ is $U(0,1)$ and independent of $J$. Hence
$R'=G+J+1$ is the rank of $T_0$ among the $B+1$ values with ties broken at random, which is
uniform on $\{1,\dots,B+1\}$ by exchangeability, and $p=(R'-1+U')/(B+1)$ with $U'$ independent of
$R'$ spreads each rank uniformly over an interval of length $1/(B+1)$. (This argument uses exact
ties; the tolerance in the code only absorbs floating-point round-off.) The time-block split (functions
\texttt{time\_block\_labels} and \texttt{train\_eval\_split}) assigns blocks using only date
ranks, block counts and independent randomness, so it satisfies (A1)(ii); for real windows it is
(A1)(i) that fails.

% ---------------------------------------------------------------------------------------------
\subsection{Consistency}\label{app:proofs-consistency}

In this subsection $P,Q$ denote the \emph{raw} window laws, $A(x)=\Lambda^{-1}(x-m_A)$ the
standardisation estimated on the pooled training half and $A^*(x)=\Lambda_*^{-1}(x-m_*)$ its
population limit. \Cref{ass:consistency} is read in standardised coordinates: (C3) says
$\ghat\in\mathcal F_n\subset\Lip_c$; (C4) says $(P-Q)(f_0\circ A^*)=\eta_0>0$; (C5) says
$(\Ptr-\Qtr)(\ghat\circ A)\ge\sup_{g\in\mathcal F_n}(\Ptr-\Qtr)(g\circ A)-\varepsilon_{\mathrm{opt}}$.
If $n_t/(n_t+m_t)$ converges and the limiting pooled law has positive coordinate (or channel)
variances, the law of large numbers gives $m_A\pto m_*$ and $\Lambda\pto\Lambda_*$ (almost surely
when the training halves are nested), hence $\norm{\Lambda^{-1}}_{\mathrm{op}}=O_\Prob(1)$.
Since $A(x)-A^*(x)=(\Lambda^{-1}-\Lambda_*^{-1})(x-m_*)+\Lambda^{-1}(m_*-m_A)$,
\begin{equation}\label{eq:app-std-error}
\begin{split}
  e_n&:=\sup_{g\in\Lip_c}\big|(P-Q)(g\circ A)-(P-Q)(g\circ A^*)\big|\\
  &\le c\big\{\norm{\Lambda^{-1}-\Lambda_*^{-1}}_{\mathrm{op}}(\E_P+\E_Q)\norm{X-m_*}
  +2\norm{\Lambda^{-1}}_{\mathrm{op}}\norm{m_A-m_*}\big\}\pto0 .
\end{split}
\end{equation}
This is the sense in which the estimation of $A$ adds $o_\Prob(1)$ terms.

\begin{lemma}[Uniform convergence over Lipschitz classes]\label{lem:uniform-lip}
Let $P\in\Pspace{1}(\R^d)$ and let $P_n$ be the empirical measure of $n$ \iid{} draws from $P$.
Then $\W(P_n,P)\to0$ almost surely and in $L^1$, and $\sup_{f\in\Lip_c}\abs{(P_n-P)f}=c\,\W(P_n,P)$.
Consequently, if the training halves have sizes $n_t,m_t\to\infty$ and are selected independently
of the data, then for every (possibly random) $c'\ge0$,
\[
  \sup_{f\in\Lip_{c'}}\big|(\Ptr-P)f-(\Qtr-Q)f\big|\le c'\big\{\W(\Ptr,P)+\W(\Qtr,Q)\big\},
\]
and $\W(\Ptr,P)+\W(\Qtr,Q)\pto0$.
\end{lemma}

\begin{proof}
By Varadarajan's theorem $P_n$ converges weakly to $P$ almost surely, and
$\int\norm{x}\,dP_n\to\int\norm{x}\,dP$ almost surely by the strong law. Together these are
equivalent to $\W(P_n,P)\to0$ \citep[Thm.~6.9]{villani2009optimal}. Coupling both measures with
$\delta_0$ gives $\W(P_n,P)\le n^{-1}\sum_i\norm{X_i}+\E\norm{X}$, a uniformly integrable
sequence, so $\E\,\W(P_n,P)\to0$ (Vitali). Every Lipschitz function is $P$- and $P_n$-integrable,
and $\Lip_c$ is symmetric under $f\mapsto-f$, so \cref{fact:kr-duality} gives the identity. Given
the index sets, the training half of $X$ is an \iid{} sample of size $n_t$, so
$\E\,\W(\Ptr,P)$ depends on $n_t$ only and tends to zero; this gives convergence in probability
along any sequence of training halves. The displayed bound holds pointwise in $\omega$.
\end{proof}

\begin{proof}[Proof of \cref{prop:consistency}]
Write $\Dcrit(f)=Pf-Qf$, $\eta=\eta_0/4$,
$\Gamma_n=c\,\norm{\Lambda^{-1}}_{\mathrm{op}}\{\W(\Ptr,P)+\W(\Qtr,Q)\}$ and
$A_n=\{\ghat\in\mathcal F_n,\ \Dcrit(\fhat)\ge\eta\}$, a $\Gsig$-measurable event. Because
$g\circ A\in\Lip_{c\norm{\Lambda^{-1}}_{\mathrm{op}}}$ for $g\in\Lip_c$, \cref{lem:uniform-lip}
bounds every training-versus-population discrepancy below by $\Gamma_n\pto0$.

\emph{Step 1: $\Prob(A_n)\to1$.} On the event of (C3)--(C5), whose probability tends to one,
\[
  \Dcrit(\fhat)\ge(\Ptr-\Qtr)(\ghat\circ A)-\Gamma_n
  \ge(\Ptr-\Qtr)(f_0\circ A)-\varepsilon_{\mathrm{opt}}-\Gamma_n
  \ge\eta_0-e_n-\varepsilon_{\mathrm{opt}}-2\Gamma_n\ge\tfrac{\eta_0}2-e_n-2\Gamma_n ,
\]
using (C4) through \eqref{eq:app-std-error}. Since $e_n,\Gamma_n\pto0$,
$\Dcrit(\fhat)\ge\eta_0/4$ with probability tending to one.

\emph{Step 2: $T$ is large.} On $A_n$, $\E[T\mid\Gsig]=\Dcrit(\fhat)\ge\eta$ by
\cref{prop:unbiased}(a), so $T\le\eta/2$ forces $\abs{T-\Dcrit(\fhat)}\ge\eta/2$, and by the
conditional Chebyshev inequality and \cref{prop:unbiased}(c),
\[
  \Prob(T\le\eta/2\mid\Gsig)\le\frac{4c^2\norm{\Lambda^{-1}}_{\mathrm{op}}^2}{\eta^2}
  \Big(\frac{\Tr\Sigma_P}{n_e}+\frac{\Tr\Sigma_Q}{m_e}\Big)=:\epsilon_{2,n}\pto0 .
\]

\emph{Step 3: permuted statistics are small.} Let $v_k=\fhat(V_k)$ on the pooled evaluation set,
with mean $\bar v$ and $S^2=N_e^{-1}\sum_k(v_k-\bar v)^2$. For a uniform permutation $\pi$, given
$(\Gsig,V)$, $t(\pi\cdot V)=(N_e/m_e)(n_e^{-1}\sum_{k\in G_1}v_k-\bar v)$ with $G_1$ a uniformly
random $n_e$-subset, and sampling without replacement gives
\[
  \E[t(\pi\cdot V)\mid\Gsig,V]=0,\qquad
  \Var(t(\pi\cdot V)\mid\Gsig,V)=\frac{N_e^2S^2}{n_em_e(N_e-1)}=:V_\pi .
\]
$S^2$ is at most the mean squared deviation about $Q\fhat$, so on $A_n$, by
\cref{prop:unbiased}(b)--(c) and \cref{rem:representation},
$\E[S^2\mid\Gsig]\le\Var_P\fhat+\Var_Q\fhat+\Dcrit(\fhat)^2\le c^2\norm{\Lambda^{-1}}_{\mathrm{op}}^2K$
with $K=\Tr\Sigma_P+\Tr\Sigma_Q+\W(P,Q)^2$.

\emph{Step 4: rejection.} Let $k_\alpha=\lfloor\alpha(B+1)\rfloor$. The test rejects iff
$\#\{b\ge1:T_b\ge T_0\}\le k_\alpha-1$. For $B\ge2/\alpha$, $k_\alpha>\alpha(B+1)-1\ge\alpha B/2\ge1$,
so $k_\alpha\ge2$ and $B/k_\alpha\le2/\alpha$. Let $M=\#\{b\ge1:T_b\ge\eta/2\}$. On
$\{T_0>\eta/2\}$, $\#\{b:T_b\ge T_0\}\le M$. Given $(\Gsig,V)$ the $T_b$ are identically
distributed, so Markov's and Chebyshev's inequalities give
$\Prob(M\ge k_\alpha\mid\Gsig,V)\le(B/k_\alpha)\,4V_\pi/\eta^2$, and on $A_n$
\[
  \Prob(M\ge k_\alpha\mid\Gsig)\le\frac{8c^2\norm{\Lambda^{-1}}_{\mathrm{op}}^2K}{\alpha\eta^2}\cdot
  \frac{N_e^2}{n_em_e(N_e-1)}=:\epsilon_{4,n}\pto0 .
\]
Hence $\Prob(\text{no rejection})\le\Prob(A_n^c)+\E[\min(1,\epsilon_{2,n})]+\E[\min(1,\epsilon_{4,n})]\to0$
by dominated convergence.
\end{proof}

\begin{proof}[Proof of \cref{cor:w1-limit}]
Let $c=1$. On $\{\ghat\in\mathcal F_n\}$, \cref{prop:unbiased}(b) and \cref{rem:representation}
give $\Dcrit(\fhat)\le\W(P_A,Q_A)$, and
$\abs{\W(P_A,Q_A)-\W(P_{A^*},Q_{A^*})}\le(\E_P+\E_Q)\norm{A(X)-A^*(X)}\pto0$ by the triangle
inequality and the coupling $(A(X),A^*(X))$. For the lower bound, the chain of Step 1 with $f_0$
replaced by near-maximisers gives
\[
  \Dcrit(\fhat)\ge\sup_{g\in\mathcal F_n}(P-Q)(g\circ A)-\varepsilon_{\mathrm{opt}}-2\Gamma_n
  \ge\sup_{g\in\mathcal F_n}(P-Q)(g\circ A^*)-e_n-\varepsilon_{\mathrm{opt}}-2\Gamma_n ,
\]
whose right-hand side tends to $\W(P_{A^*},Q_{A^*})$ in probability by assumption. Finally
$T-\Dcrit(\fhat)\pto0$ by Step 2 of the previous proof (Chebyshev with any $\eta>0$).
\end{proof}

% ---------------------------------------------------------------------------------------------
\subsection{Likelihood connection and the Neyman--Pearson bound}\label{app:proofs-likelihood}

\begin{proof}[Proof of \cref{prop:oracle-bound}]
(a) Let $\bar\phi(x,y)=\int\phi(x,y,u)\,\Prob_U(du)$. Because $U$ is independent of the data,
Fubini's theorem gives $\E_F\phi=\E_F\bar\phi$ under every data law $F$, so $\bar\phi$ is a
randomised test with the same level and power. Consider the simple problem
$H_0':(X,Y)\sim P^{\otimes n}\otimes P^{\otimes m}$ against
$H_1':(X,Y)\sim P^{\otimes n}\otimes Q^{\otimes m}$. The likelihood ratio is $\prod_jq(Y_j)/p(Y_j)$
(the $X$-factors cancel), so the most powerful level-$\alpha$ test, given by the
Neyman--Pearson lemma \citep[Thm.~3.2.1]{lehmann2022testing}, depends on $Y$ only and has power
$\bNP(P\to Q;m,\alpha)$. $\bar\phi$ has level $\alpha$ for $H_0'$, so its power under $H_1'$ is at
most that. (b) is the same argument with the null point $Q^{\otimes n}\otimes Q^{\otimes m}$; the
likelihood ratio is then $\prod_ip(X_i)/q(X_i)$.
\end{proof}

\begin{proof}[Proof of \cref{prop:local-power}]
Let $\Lambda_m=\sum_{j\le m}\log\{q_{\theta_m}(Y_j)/p(Y_j)\}$ be the log-likelihood ratio of
$P^{\otimes n}\otimes Q_{\theta_m}^{\otimes m}$ against $P^{\otimes n}\otimes P^{\otimes m}$. By
differentiability in quadratic mean, under $\theta=0$,
$\Lambda_m=\tau m^{-1/2}\sum_{j\le m}h(Y_j)-\tau^2I/2+o_\Prob(1)$ \citep[Thm.~7.2]{vandervaart_1998_asymptotic},
so $\Lambda_m\dto N(-\tau^2I/2,\tau^2I)$ and the alternatives are contiguous.

(a) The most powerful test rejects for large $\Lambda_m$. By Le Cam's third lemma
\citep[Ex.~6.7]{vandervaart_1998_asymptotic}, $\Lambda_m\dto N(\tau^2I/2,\tau^2I)$ under $\theta_m$.
The limit under the null is continuous, so the exact level-$\alpha$ critical values converge to
$-\tau^2I/2+z_{1-\alpha}\tau\sqrt I$ and the power converges to $1-\Phi(z_{1-\alpha}-\tau\sqrt I)$.

(b) The evaluation index sets are chosen independently of the data and $f$ is fixed. Under
$\theta=0$,
$\sqrt mT=\sqrt m\,n_e^{-1}\sum_{i\in\Iev}\{f(X_i)-Pf\}-\sqrt m\,m_e^{-1}\sum_{j\in\Jev}\{f(Y_j)-Pf\}$.
The multivariate Lindeberg central limit theorem applied to $(\sqrt mT,\Lambda_m)$, whose $X$-part
is independent of the $Y$-part, gives joint asymptotic normality with
\[
  m\Var T\to\sigma_f^2\Big(\frac1{\rho\kappa}+\frac1\rho\Big)=\sigma_f^2\,\frac{1+\kappa}{\rho\kappa},
  \qquad
  \Cov(\sqrt mT,\Lambda_m)=-\frac{\tau}{m_e}\sum_{j\in\Jev}\Cov_P(f,h)\to\tau r_f\sigma_f\sqrt I ,
\]
since only the $m_e$ synthetic evaluation windows contribute to the covariance and
$-\Cov_P(f,h)=r_f\sigma_f\sqrt I$. By Le Cam's third lemma, under $\theta_m$,
$\sqrt mT\dto N\big(\tau r_f\sigma_f\sqrt I,\ \sigma_f^2(1+\kappa)/(\rho\kappa)\big)$. The variance
estimators $s_X^2,s_Y^2$ converge to $\sigma_f^2$ under $\theta=0$, hence under $\theta_m$ by
contiguity (Le Cam's first lemma), so $m(s_X^2/n_e+s_Y^2/m_e)\pto\sigma_f^2(1+\kappa)/(\rho\kappa)$
and $z\dto N(\tau\sqrt I\,r_f\sqrt{\rho\kappa/(1+\kappa)},1)$, which gives the stated power. For the
permutation test, under $\theta=0$ the pooled evaluation values are \iid, and the permutation
distribution of $\sqrt mT_b$ converges conditionally in probability to
$N(0,\sigma_f^2(1+\kappa)/(\rho\kappa))$ \citep{hoeffding_1952_large}; its $(1-\alpha)$ quantile
converges in probability to $z_{1-\alpha}\sigma_f\{(1+\kappa)/(\rho\kappa)\}^{1/2}$, also under
$\theta_m$ by contiguity, and the limit power is the same. This holds for the test based on the
permutation quantile, that is with $B\to\infty$; with a fixed number $B$ of random permutations the
Monte Carlo p-value has a smaller limiting power.

(c) $r_f\le1$ by the Cauchy--Schwarz inequality, and $\sqrt{\rho\kappa/(1+\kappa)}=1/2$ at
$\kappa=1$, $\rho=1/2$. Where the oracle has limiting power $0.80$ at $\alpha=0.05$, its
noncentrality is $\tau\sqrt I=z_{0.95}+z_{0.80}\approx2.486$, so every test based on the held-out
mean difference of a fixed critic has noncentrality at most $1.243$ and limiting power at most $1-\Phi(1.645-1.243)\approx0.34$.

\emph{Learned critics (sketch).} Conditionally on the training half, $\fhat$ is a fixed function
of the training data, and the evaluation data, independent of it, play the role of the data in
(b); the bound $r_{\fhat}\le1$ then holds conditionally. Making the argument rigorous along a
sequence in which $\fhat$ changes with $m$ needs the joint convergence in (b) to hold uniformly
over the class that contains $\fhat$ (for instance $\Lip_c$ with $P\in\Pspace{2+\epsilon}$),
which we do not prove.
\end{proof}

\begin{proof}[Details of \cref{ex:vol-level}]
With $q_\theta(x)\propto(1+\theta)^{-L}\exp\{-\norm{x}^2/(2(1+\theta)^2\sigma^2)\}$, the score at
$\theta=0$ is $h(x)=\norm{x}^2/\sigma^2-L$. For $s=1+\theta>1$, the coupling $Y=sX$ gives
$\W(P,Q_\theta)\le(s-1)\E\norm X$, and the 1-Lipschitz function $f(x)=-\norm x$ attains it:
$Pf-Q_\theta f=(s-1)\E\norm X$. So $f$ is an optimal Kantorovich potential that scores real
windows higher, for every $\theta>0$. Its orientation is
$r_f=-\Corr_P(f,h)=\Corr(\chi_L,\chi_L^2)$, where $\chi_L=\norm X/\sigma$. With
$\E\chi_L^j=2^{j/2}\Gamma((L+j)/2)/\Gamma(L/2)$, $\Var\chi_L^2=2L$ and
$\Cov(\chi_L,\chi_L^2)=\E\chi_L^3-L\,\E\chi_L$, one obtains $r_f^2=0.876$, $0.957$ and $0.988$ for
$L=1,5,20$. The pooled law $\tfrac12(P+Q_\theta)$ has mean zero and equal coordinate variances, so
the population standardisation is $A^*(x)=x/c$ for a constant $c>0$; it rescales $f$, and a
correlation is scale invariant.
\end{proof}

% ---------------------------------------------------------------------------------------------
\subsection{The relative tests}\label{app:proofs-relative}

Here $\Gsig$ is generated by $X$, $Y^{\trn}$, $Z^{\trn}$, the index sets and $\xi$, as in
\cref{sec:theory-relative}. The evaluation vector is $V=(Z^{\evl},Y^{\evl})$, past windows first,
and $t(v)=k_e^{-1}\sum_{l\le k_e}\fhat(v_l)-m_e^{-1}\sum_{l>k_e}\fhat(v_l)$, so that $\Danc=t(V)$.
This is how \texttt{relative\_test\_anchored} calls the permutation routine (values of $\fhat$ on
$Z^{\evl}$ first, then on $Y^{\evl}$; $\fhat$ frozen; $\Rnew$ never evaluated).

\begin{proof}[Proof of \cref{prop:anchored}]
(a) By \cref{ass:anc}(i), conditionally on $X$ the $k+m$ windows of $Z$ and $Y$ are \iid. By (ii)
the index sets are functions of time stamps, group labels and randomness independent of the
window values, so, conditionally on $X$ and the index sets, the evaluation windows are \iid{} and
independent of the training windows. By (iii) $\fhat$ is a function of $X$, $Y^{\trn}$, $Z^{\trn}$
(through the standardisation as well as through training) and of $\xi$, which is independent of
the data. Hence, conditionally on $\Gsig$, the $k_e+m_e$ coordinates of $V$ are \iid{}, and $V$ is
exchangeable: (A1$'$) holds for $V$ with the statistic $t$. With (iv), \cref{lem:exchangeable,lem:rank}
give $\Prob(\pperm\le\alpha\mid\Gsig)\le\lfloor\alpha(B+1)\rfloor/(B+1)$ exactly as in the proof of
\cref{prop:exact-level}. Nothing in this argument involves $\Pnew$ or the training objective: the
pool weights, the gradient penalty and early stopping affect only $\fhat$, which is
$\Gsig$-measurable. The tolerance in the tie count only increases $\pperm$.

(b) Conditionally on $\Gsig$, $\fhat$ is fixed, and the problem is a two-sample comparison of the
means of $\fhat(Z)$, $Z\sim\Ppast$, and $\fhat(Y)$, $Y\sim\QG$, with independent \iid{} samples of
sizes $k_e,m_e$ and $k_e/(k_e+m_e)\to\kappa$. Write $\sigma_Z^2=\Var_{\Ppast}\fhat$,
$\sigma_Y^2=\Var_{\QG}\fhat$ and $N=k_e+m_e$. For the studentised statistic
$\Danc/(s_Z^2/k_e+s_Y^2/m_e)^{1/2}$, the permutation distribution converges conditionally in
probability to $N(0,1)$ whatever the two laws, provided the variances are finite and positive
\citep{janssen1997studentized,chung2013exact}. The observed statistic is asymptotically $N(0,1)$ on
$\{\delta=0\}$ (Lindeberg) and tends to $-\infty$ in probability on $\{\delta<0\}$, because its
numerator converges to $\delta$ and its denominator to zero. Hence
$\limsup\Prob(\text{reject}\mid\Gsig)\le\alpha$ on $\{\delta\le0\}$. For the unstudentised
statistic, the permutation distribution of $\sqrt N\Danc_b$ converges to $N(0,\bar\sigma^2/\{\kappa(1-\kappa)\})$,
where $\bar\sigma^2$ is the variance of $\fhat$ under the mixture $\kappa\Ppast+(1-\kappa)\QG$
\citep{chung2013exact}. On $\{\delta=0\}$ the means agree, $\bar\sigma^2=\kappa\sigma_Z^2+(1-\kappa)\sigma_Y^2$,
and the permutation variance $\sigma_Z^2/(1-\kappa)+\sigma_Y^2/\kappa$ equals the true asymptotic
variance $\sigma_Z^2/\kappa+\sigma_Y^2/(1-\kappa)$ of $\sqrt N\Danc$ when $\kappa=1/2$. On
$\{\delta<0\}$, $\sqrt N\Danc\to-\infty$ while the permutation quantile stays bounded in
probability. So for $k_e=m_e$ the unstudentised test is asymptotically valid on $\{\delta\le0\}$.
When $k_e\neq m_e$ and $\sigma_Z\neq\sigma_Y$ it is not, which is why the code also reports the
studentised permutation p-value (\texttt{p\_perm\_stud}).

(c) Conditionally on $\Gsig$, \cref{ass:mixing} holds for the fixed function $f=\fhat$, with the
evaluation windows of $Z$ in the role of the real sample and those of $Y$ in the role of the
synthetic one. The proof of \cref{prop:hac} below shows, without using $H_0'$, that
$(\Danc-\delta)/\{\widehat V_{\mathrm{DK}}(\bar a)+\widehat V_{\mathrm{DK}}(\bar b)\}^{1/2}\dto N(0,1)$, where
$\bar a=Z^{\evl}\fhat$, $\bar b=Y^{\evl}\fhat$. On $\{\delta=0\}$ this is $\zhac\dto N(0,1)$. On
$\{\delta<0\}$ the denominator tends to zero in probability (it is ratio-consistent for
$\Var(\bar a)+\Var(\bar b)\to0$), so $\zhac=(\Danc-\delta)/\hat s+\delta/\hat s\to-\infty$. The
test that rejects when $\zhac>z_{1-\alpha}$ therefore has conditional asymptotic level $\alpha$ on
$\{\delta\le0\}$, and by dominated convergence
$\limsup\Prob(\text{reject},\,\delta\le0)\le\alpha\,\Prob(\delta\le0)$.
\end{proof}

\emph{What the proof does not cover in the implemented protocol.} Part (a) concerns the
permutation p-value and needs the $Z$ and $Y$ windows to be \iid{} given $X$, as in E2 part e.
The shared block labels of the time-block split (a tenth of the \TEST{} horizon of $G$ and the
corresponding tenth of \TRAIN{} fall in the same half) are allowed by \cref{ass:anc}(ii). The
verdict rule combines the Driscoll--Kraay p-values of part (c), which are valid only in the limit,
also for independent windows (\cref{prop:verdict}(a)). Part (c) needs evaluation windows
independent of $\Gsig$, which the interleaved split without a gap does not deliver (see after the
proof of \cref{lem:decoupling}), and a stationary past. On real data neither holds, $G$ is
calibrated on $\Rpast$, and $\delta$ is random and not centred at zero (\cref{rem:anchored-null},
\cref{rem:nonstat-past}). The verdict therefore relies on the estimated boundary size of the
seed-combined rule in designs D1 and D2 and on the approximate calibration of
\cref{rem:calibrated-bad}, not on \cref{prop:anchored}.

\begin{proof}[Proof of \cref{prop:anchored-target}]
Write $\fhat=\ghat\circ A$, so $\delta=(A_\#\Ppast)\ghat-(A_\#\QG)\ghat$.
(a) For $c=0$ there is nothing to prove. For $c>0$, $\pm\ghat/c\in\Lip_1$ is integrable under
both standardised laws, and \cref{fact:kr-duality} gives $\pm\delta/c\le\W(A_\#\Ppast,A_\#\QG)$. If
$\QG=\Ppast$ the two means coincide for any integrable $\fhat$, whatever $\Pnew$.

(b) With equal class weights, the population problem is \cref{prop:bayes-logit} with $\pi=1/2$,
positive class $\Pnew$ and negative class $M$. Mutual absolute continuity of $\Pnew,\QG,\Ppast$
implies that of $\Pnew$ and $M$, so $\lstar=\log(d\Pnew/dM)$ is finite almost everywhere. Since
$M\ge\QG/2$, $d\QG/dM\in(0,2]$, and $\QG$-a.s.
\[
  \lstar=\log\frac{d\Pnew}{dM}=-\log\frac{d\QG}{d\Pnew}+\log\frac{d\QG}{dM}.
\]
Both terms are $\QG$-integrable: $\KL(\QG\Vert\Pnew)<\infty$ controls the positive part of the
first logarithm, and for any $R\ll S$ the negative part of $\log(dR/dS)$ has $R$-integral
$\E_R\log^+(dS/dR)\le\E_R(dS/dR)\le1$; the second logarithm is bounded above by $\log2$. Integrating,
$\QG\lstar=-\KL(\QG\Vert\Pnew)+\KL(\QG\Vert M)$, and likewise
$\Ppast\lstar=-\KL(\Ppast\Vert\Pnew)+\KL(\Ppast\Vert M)$; subtracting gives \eqref{eq:anchored-kl}.
Because $d\QG/dM\le2$, $0\le\KL(\QG\Vert M)\le\log2$, and the same holds for $\Ppast$; so the
second bracket lies in $[-\log2,\log2]$. For the expansion, take $\varepsilon\norm{h}_\infty<1$
and write the densities with respect to $\Ppast$ as $1$, $1+x$ and $1+x/2$ with $x=\varepsilon h$.
Then $\KL(\QG\Vert M)-\KL(\Ppast\Vert M)=\Ppast\upsilon(\varepsilon h)$, where
\[
  \upsilon(x)=(1+x)\log(1+x)-x\log(1+x/2)=x-\tfrac{x^3}{24}+\tfrac{x^4}{24}+O(x^5),
\]
with a remainder bounded by $C\abs x^5$ uniformly on $\abs x\le\varepsilon\norm h_\infty$. Since
$\Ppast h=0$, the second bracket equals $-\varepsilon^3\Ppast(h^3)/24+O(\varepsilon^4)$.

(c) $\delta=\Ppast f^*-\QG f^*=(\Pnew-\QG)f^*-(\Pnew-\Ppast)f^*$ is an identity. The sum of the two
terms is $2\Pnew f^*-(\QG+\Ppast)f^*=2(\Pnew-M)f^*$, and since $f^*=g^*\circ A$ with $g^*$
1-Lipschitz in standardised coordinates maximises $(\Pnew-M)f$ over this class,
\cref{fact:kr-duality} gives $(\Pnew-M)f^*=\W(A_\#\Pnew,A_\#M)$.
\end{proof}

\begin{proof}[Proof sketch of \cref{prop:paired-boot}]
Conditionally on $\Gsig$ both critics are fixed. (a) $\Dtwo-\delta^{\mathrm{2c}}$ is the sum of
three independent centred sample means, of $\fhat_G-\fhat_P$ over $X^{\evl}$, of $-\fhat_G$ over
$Y^{\evl}$ and of $\fhat_P$ over $Z^{\evl}$, with variance $\sigma_N^2$; the Lyapunov condition
of (R2) gives the central limit theorem. (b) Each bootstrap mean satisfies a conditional central
limit theorem with the empirical variance \citep[Thm.~23.4]{vandervaart_1998_asymptotic}; for
triangular arrays the Lyapunov condition for the bootstrap draws holds in probability by the law
of large numbers for the empirical $(2+\epsilon)$-th moments. The three bootstrap means are
independent given the data, and P\'olya's theorem makes the convergence uniform because the limit
is continuous. Pairing makes the $X$-term the bootstrap mean of the differences
$\fhat_G(X_i)-\fhat_P(X_i)$, whose bootstrap variance estimates $\sigma_X^2$ including the
covariance of the two critics. (c) By (b), the bootstrap $(1-\alpha)$ quantile $q^*$ of
$\Dtwo{}^*-\Dtwo$ satisfies $q^*/\sigma_N\pto z_{1-\alpha}$ as $B\to\infty$, and
$\Prob(\Dtwo>q^*\mid\Gsig)=\Prob\{(\Dtwo-\delta^{\mathrm{2c}})/\sigma_N>q^*/\sigma_N-\delta^{\mathrm{2c}}/\sigma_N\mid\Gsig\}$;
apply (a). The centred p-value $p^*$ rejects exactly when $\Dtwo$ exceeds this quantile
\citep{hall1991two}. (d) The percentile interval contains $\delta^{\mathrm{2c}}$ iff
$(\Dtwo-\delta^{\mathrm{2c}})/\sigma_N\in[-z_{0.975},z_{0.975}]+o_\Prob(1)$, an event of probability
$0.95+o(1)$ by (a), using the symmetry of the normal limit.
\end{proof}

\begin{proof}[Proof of \cref{prop:boundary-size}]
(a) By \cref{prop:paired-boot}(c),
$\Prob(\text{reject}\mid\Gsig)-\{1-\Phi(z-\delta^{\mathrm{2c}}/\sigma_N)\}\pto0$; the difference
is bounded by one, so its expectation tends to zero.
(b) Let $\varpi$ collect $X,Y,Z,\xi_G,\xi_P$ and the index sets, and let $\varpi'$ be $\varpi$ with $(Y,\xi_G)$
and $(Z,\xi_P)$, and their index sets, exchanged; by assumption $\varpi'\eqd\varpi$. The exchange
maps $\fhat_G=\Psi(X^{\trn},Y^{\trn},\xi_G)$ to $\fhat_P$ and conversely. Because $\QG=\Ppast$, it
exchanges $D_G=(\Pnew-\QG)\fhat_G$ and $D_P=(\Pnew-\Ppast)\fhat_P$; it leaves
$\sigma_X^2=\Var_{\Pnew}(\fhat_G-\fhat_P)$ unchanged; and, using $m_e=k_e$, it exchanges
$\sigma_Y^2/m_e$ and $\sigma_Z^2/k_e$. Hence
$(\delta^{\mathrm{2c}},\sigma_N)(\varpi')=(-\delta^{\mathrm{2c}},\sigma_N)(\varpi)$, which gives the
symmetry. For a symmetric $U=\delta^{\mathrm{2c}}/\sigma_N$,
\[
  \E\{1-\Phi(z-U)\}=\E\,\tfrac12\big[\{1-\Phi(z-\abs U)\}+\{1-\Phi(z+\abs U)\}\big]=\E\,\bar g(\abs U).
\]
$\bar g$ is continuous with $\bar g(0)=1-\Phi(z)=\alpha$ and $\bar g(x)\to1/2$ as $x\to\infty$, and
$\bar g'(x)=\tfrac12\{\varphi(z-x)-\varphi(z+x)\}>0$ for $x>0$ because $z>0$ gives
$\abs{z-x}<z+x$. The two bounds follow from $\bar g\ge\alpha$ on $[0,\infty)$ and
$\bar g\ge\bar g(\varepsilon)$ on $[\varepsilon,\infty)$.
(c) Now $\delta^{\mathrm{2c}}$ does not change with the evaluation sizes and $\sigma_N\to0$, so
$1-\Phi(z-\delta^{\mathrm{2c}}/\sigma_N)\to\one\{\delta^{\mathrm{2c}}>0\}+\alpha\one\{\delta^{\mathrm{2c}}=0\}$;
apply (a) and dominated convergence. Under the symmetry of (b),
$\Prob(\delta^{\mathrm{2c}}>0)=\Prob(\delta^{\mathrm{2c}}<0)$, which gives $1/2$ without an atom at
zero.
\end{proof}

% ---------------------------------------------------------------------------------------------
\subsection{Serial and cross-sectional dependence}\label{app:proofs-dependence}

\begin{proof}[Proof of \cref{prop:hac}]
Let $c_X$ and $c_Y$ be the fixed numbers of windows per date and $\mathcal D_X$, $\mathcal D_Y$ the
sets of evaluation dates, with $\abs{\mathcal D_X},\abs{\mathcal D_Y}\to\infty$ proportionally by
(D4). Then $\bar a=(c_X\abs{\mathcal D_X})^{-1}\sum_{\tau\in\mathcal D_X}u_\tau$ and
$\E u_\tau=c_XPf$. First let $\mathcal D_X$ be a set of consecutive dates. Under (D1)--(D2), the
central limit theorem for strongly mixing stationary sequences \citep{herrndorf1984functional},
whose conditions are exactly the moment and mixing conditions of (D1) with $\beta=2+\epsilon$,
gives $\abs{\mathcal D_X}^{1/2}(\bar a-Pf)\dto N(0,\LRV_u/c_X^2)$ with $\LRV_u>0$, and
$\abs{\mathcal D_X}\Var(\bar a)\to\LRV_u/c_X^2$ because the covariances of $(u_\tau)$ are summable
under (D1). The same holds for $\bar b$, and the two sample means are independent. Let
$s_N^2=\Var(\bar a)+\Var(\bar b)$. Then $\{(\bar a-Pf)-(\bar b-Qf)\}/s_N\dto N(0,1)$, because along
every subsequence on which $\Var(\bar a)/s_N^2$ converges the limit is a sum of two independent
centred normals with variances summing to one. By (D3),
$\{\widehat V_{\mathrm{DK}}(\bar a)+\widehat V_{\mathrm{DK}}(\bar b)\}/s_N^2\pto1$, and Slutsky's lemma gives
$\{(\bar a-\bar b)-(Pf-Qf)\}/\{\widehat V_{\mathrm{DK}}(\bar a)+\widehat V_{\mathrm{DK}}(\bar b)\}^{1/2}\dto N(0,1)$.
Under $H_0'$ this is $\zhac\dto N(0,1)$.
\end{proof}

\emph{Evaluation dates that form a union of time blocks.} Under the time-block split, the
evaluation dates of a sample are the dates whose block label lies in a set $\mathcal E$ of labels,
drawn independently of the data; we condition on $\mathcal E$. The code
(\texttt{time\_block\_labels}) cuts the sample's date range into $K$ intervals of equal length, so
block $j$ consists of the dates whose relative position lies in $[j/K,(j+1)/K)$. Let
$\mathbb S_T(r)=T^{-1/2}\sum_{\tau\le\tau_{\min}+rT}(u_\tau-c_XPf)$ be the partial-sum process
over all $T$ dates of the range. \citet{herrndorf1984functional} proves, under the same conditions,
the functional limit $\mathbb S_T\dto\LRV_u^{1/2}\mathbb B$ with $\mathbb B$ a standard Brownian
motion. The evaluation sum equals
$T^{1/2}\sum_{j\in\mathcal E}\{\mathbb S_T((j+1)/K)-\mathbb S_T(j/K)\}$ up to at most $2K$ boundary
terms, each $O_\Prob(1)$. By the continuous mapping theorem it is asymptotically normal with
variance $T\LRV_u\abs{\mathcal E}/K$, and $\abs{\mathcal D_X}/T\to\abs{\mathcal E}/K$; so the
central limit theorem above holds for the union. For the variance estimator, the code places the
evaluation windows on the full grid of window dates and counts the training positions as zeros
(\texttt{hac\_var\_mean}, \eqref{eq:dk-var}). Products $\mathcal S_u\mathcal S_{u+h}$ with one
position in a training block vanish, so $\widehat V_{\mathrm{DK}}(\bar a)$ is the sum of the Bartlett
estimators of the maximal runs of consecutive evaluation dates plus products that straddle a
training block shorter than $H$. There are at most $K$ such gaps and at most $H(H+1)/2$ straddling
products per gap, each with bounded expectation. If the Bartlett estimator is consistent on each
run and $KH^2/T\to0$, which holds for the Andrews plug-in rate $H\asymp T^{1/3}$ with $K=10$ fixed,
(D3) therefore carries over to the union. The Andrews AR(1) coefficient is computed over adjacent
observed positions only, so the holes do not bias the choice of $H$ towards zero. The argument
treats $f$ as fixed. For the learned critic it applies conditionally on $\Gsig$ only when the
evaluation windows are independent of $\Gsig$, which the interleaved split does not guarantee.

\begin{proof}[Proof of \cref{lem:decoupling}]
Regular conditional distributions exist on Polish spaces; let $\Prob_{V\mid U}$ denote that of $V$
given $U$, and $E_u=\{v:(u,v)\in E\}$. Then $\Prob((U,V)\in E)=\E\,\Prob_{V\mid U}(E_U)$ and, since
$\tilde V$ is independent of $U$ with the law $\Prob_V$ of $V$,
$\Prob((U,\tilde V)\in E)=\E\,\Prob_V(E_U)$. Each $E_U$ is a measurable set of the range of $V$, so
the difference is at most
$\E\sup_C\abs{\Prob_{V\mid U}(C)-\Prob_V(C)}=\E\sup_{B\in\mathcal B}\abs{\Prob(B\mid\mathcal A)-\Prob(B)}=\beta(\mathcal A,\mathcal B)$.
\end{proof}

The bound $\beta(\mathcal A,\mathcal B)\le\beta_r(g+1)$ for contiguous blocks separated by a
gap of $g$ days uses two properties of $\beta$: it is monotone in each $\sigma$-field, and
adjoining components independent of everything else (training randomness, independent synthetic
windows) does not increase it \citep[Thm.~6.2]{bradley_2005_basic}. Under the implemented split
($K=10$ interleaved blocks, no gap) the training and evaluation $\sigma$-fields contain those of
two adjacent blocks, so by monotonicity $\beta(\mathcal A,\mathcal B)$ is at least a lag-one
coefficient of the panel, which is not small under persistent volatility. The lemma then gives no
useful bound, \cref{cor:block-bias} shows how a persistent common factor can then shift the mean
of $T$, and the residual dependence is assessed by simulation (E1, E2 part g).

\begin{proof}[Proof of \cref{prop:ticker-bias}]
(a) The real evaluation windows are functions of $F$ and of $(\varepsilon_{i\tau})_{i\in\Iev}$.
$\Gsig$ depends on the real data only through $F$ and $(\varepsilon_{i\tau})_{i\in\Itr}$, and on
synthetic training windows, the split and $\xi$, all independent of
$(\varepsilon_{i\tau})_{i\in\Iev}$. Hence, given $(\Gsig,F)$, the windows $X_{i\tau}$, $i\in\Iev$,
are independent with laws $P^{(F_\tau)}$, and $\fhat$ is fixed given $(\Gsig,F)$. The synthetic
evaluation windows are independent of $(\Gsig,F)$ with law $\bar Q$. Each evaluation ticker contributes one window per date of
$\mathcal T$, so, for $\fhat$ measurable with respect to $\Gsig\vee\sigma(F)$ and integrable under
every $P^{(v)}$ and under $\bar Q$,
$\E[\Pev\fhat\mid\Gsig,F]=\abs{\mathcal T}^{-1}\sum_{\tau\in\mathcal T}P^{(F_\tau)}\fhat=\hat P_F\fhat$ and
$\E[\Qev\fhat\mid\Gsig,F]=\bar Q\fhat$.
(b) Take expectations in (a). If $\fhat$ is independent of $F$, condition on $\fhat$ and use
$\E\,P^{(F_\tau)}g=\int P^{(v)}g\,\mu_F(dv)=\bar Pg=\bar Qg$ for every fixed integrable $g$ and every
$\tau$, by stationarity of $(F_\tau)$.
(c) The idealised critic is a function of $F$, so (a) applies to it. If $\hat P_F$ and $\bar Q$
are mutually absolutely continuous, (a) gives
$\E[T\mid\Gsig,F]=\hat P_F\log(d\hat P_F/d\bar Q)-\bar Q\log(d\hat P_F/d\bar Q)
=\KL(\hat P_F\Vert\bar Q)+\KL(\bar Q\Vert\hat P_F)=\Jeff(\hat P_F,\bar Q)$, which is positive,
possibly infinite, when the laws differ. In the Gaussian scale mixture both laws have positive densities. The characteristic
function of a mixture $\int P^{(v)}\varrho(dv)$ at $t\in\R^L$ is
$\int e^{-v^2\norm t^2/2}\varrho(dv)$, the Laplace transform at $\norm t^2/2$ of the law $\varrho_2$
of $v^2$ under $\varrho$, so the mixture determines $\varrho_2$. For $\hat P_F$,
$\varrho=\abs{\mathcal T}^{-1}\sum_\tau\delta_{F_\tau}$ and $\varrho_2$ is purely atomic; for
$\bar Q=\bar P$, $\varrho=\mu_F$, and $\varrho_2$ has no atoms when $\mu_F$ has none, because
$v\mapsto v^2$ is injective on $(0,\infty)$. Hence $\hat P_F\neq\bar Q$ almost surely.
\end{proof}

% ---------------------------------------------------------------------------------------------
\subsection{Window distances and portfolio utility}\label{app:proofs-markowitz}

Here $\Sigma$ is the \TEST{} covariance matrix (population, or the realised $S_{\mathrm{te}}$, in
which case every identity holds exactly for the measured quantities), $\mu$ the \TEST{} mean,
$V(w)=w'\Sigma w$, $w^*=\wgmv(\Sigma)$, $V^*=V(w^*)$, $c_{\one}=\one'\Sigma^{-1}\one$, and
$M=\Sigma^{-1}-\Sigma^{-1}\one\one'\Sigma^{-1}/c_{\one}$. For $\hat\Sigma$ positive definite, let
$A_\Sigma=\Sigma^{-1/2}\hat\Sigma\Sigma^{-1/2}$, which is similar to $\Sigma^{-1}\hat\Sigma$; the
condition number $\kappa=\cond(\Sigma^{-1}\hat\Sigma)$ of \eqref{eq:kantorovich-preview} is read as
$\lambda_{\max}(A_\Sigma)/\lambda_{\min}(A_\Sigma)$, the ratio of the extreme (real, positive)
eigenvalues of $\Sigma^{-1}\hat\Sigma$, which differs in general from the spectral condition number
$\norm{\Sigma^{-1}\hat\Sigma}_{\mathrm{op}}\norm{\hat\Sigma^{-1}\Sigma}_{\mathrm{op}}$ of the
non-symmetric product.

\begin{proof}[Proof of \cref{lem:quad-loss}]
(a) $\Sigma w^*=\one/c_{\one}$, so $w^{*\prime}\Sigma(w-w^*)=\one'(w-w^*)/c_{\one}=0$ whenever $\one'w=1$, and
$V(w)=V^*+2w^{*\prime}\Sigma(w-w^*)+(w-w^*)'\Sigma(w-w^*)$.
(b) For a closed convex feasible set, first-order optimality of $w^*$ gives
$w^{*\prime}\Sigma(w-w^*)\ge0$ for feasible $w$; use the same expansion.
(c) With the budget constraint only, the first-order condition is $\mu-\gamma\Sigma\wmv^*=\nu\one$
for a multiplier $\nu$, and for feasible $w$
$\CE(w)=\CE(\wmv^*)+(\mu-\gamma\Sigma\wmv^*)'(w-\wmv^*)-\tfrac\gamma2(w-\wmv^*)'\Sigma(w-\wmv^*)$,
whose middle term is $\nu\one'(w-\wmv^*)=0$. With further convex constraints the middle term is
$\le0$ by first-order optimality.
\end{proof}

\begin{lemma}[Kantorovich inequality]\label{lem:kantorovich-ineq}
Let $\mathsf B$ be symmetric positive definite with spectrum in $[\underline m,\overline m]$,
$0<\underline m<\overline m$. For every $x$,
$(x'\mathsf Bx)(x'\mathsf B^{-1}x)\le\frac{(\underline m+\overline m)^2}{4\underline m\overline m}(x'x)^2$,
with equality if and only if $x=\alpha e_{\underline m}+\beta e_{\overline m}$ with $\alpha^2=\beta^2$,
where $e_{\underline m},e_{\overline m}$ are unit eigenvectors of the extreme eigenvalues.
\end{lemma}

\begin{proof}
$(\mathsf B-\underline mI)(\overline mI-\mathsf B)\succeq0$ gives
$\mathsf B+\underline m\overline m\,\mathsf B^{-1}\preceq(\underline m+\overline m)I$, hence
$2\{\underline m\overline m(x'\mathsf Bx)(x'\mathsf B^{-1}x)\}^{1/2}\le x'\mathsf Bx+\underline m\overline m\,x'\mathsf B^{-1}x\le(\underline m+\overline m)x'x$
by the inequality of arithmetic and geometric means; square. Equality in the second step requires
$x$ to lie in the sum of the two extreme eigenspaces, and equality in the first requires
$x'\mathsf Bx=\underline m\overline m\,x'\mathsf B^{-1}x$; for $x=\alpha e_{\underline m}+\beta e_{\overline m}$
this reads $\underline m\alpha^2+\overline m\beta^2=\overline m\alpha^2+\underline m\beta^2$, that is
$\alpha^2=\beta^2$.
\end{proof}

\begin{proof}[Proof of \cref{prop:gmv-sensitivity}]
(a) $\wgmv(c\hat\Sigma)=c^{-1}\hat\Sigma^{-1}\one/(c^{-1}\one'\hat\Sigma^{-1}\one)=\wgmv(\hat\Sigma)$.
(b) Let $u=\Sigma^{-1/2}\one$. Since $\hat\Sigma^{-1}=\Sigma^{-1/2}A_\Sigma^{-1}\Sigma^{-1/2}$,
$\one'\hat\Sigma^{-1}\one=u'A_\Sigma^{-1}u$, $\one'\hat\Sigma^{-1}\Sigma\hat\Sigma^{-1}\one=u'A_\Sigma^{-2}u$
and $V^*=1/(u'u)$, so
$V(\wgmv(\hat\Sigma))/V^*=(u'A_\Sigma^{-2}u)(u'u)/(u'A_\Sigma^{-1}u)^2$. Apply
\cref{lem:kantorovich-ineq} with $\mathsf B=A_\Sigma^{-1}$ and $x=\mathsf B^{1/2}u$; the spectrum
of $\mathsf B$ has extreme ratio $\kappa$. The lower bound is the optimality of $w^*$. For
sharpness, let $N\ge2$ and $\kappa\ge1$, choose orthonormal $e_1,e_2$ spanning a plane that
contains $u$, write $u=c_1e_1+c_2e_2$ with $c_2^2/c_1^2=\kappa$, and let $A_\Sigma$ have
eigenvalue $1$ on $e_1$, $\kappa$ on $e_2$ and any eigenvalues in $[1,\kappa]$ on the orthogonal
complement. Then $A_\Sigma^{-1/2}u=c_1e_1+c_2\kappa^{-1/2}e_2$ has equal mass on the two extreme
eigenvectors, equality holds, and $\hat\Sigma=\Sigma^{1/2}A_\Sigma\Sigma^{1/2}$ attains the bound.
(c) $\norm{cA_\Sigma-I}_{\mathrm{op}}=\max\{\abs{c\lambda_{\min}-1},\abs{c\lambda_{\max}-1}\}$ is
minimised at $c=2/(\lambda_{\max}+\lambda_{\min})$, with value $(\kappa-1)/(\kappa+1)$, and
$\varepsilon^2/(1-\varepsilon^2)=(\kappa-1)^2/(4\kappa)=(1+\kappa)^2/(4\kappa)-1$. If
$\norm{A_\Sigma-I}_{\mathrm{op}}=\varepsilon_0<1$, then $\varepsilon\le\varepsilon_0$ (take $c=1$),
and $x\mapsto x^2/(1-x^2)$ is increasing on $[0,1)$.
(d) $(\Sigma+E)^{-1}=\Sigma^{-1}-\Sigma^{-1}E\Sigma^{-1}+O(\norm E^2)$, $\Sigma^{-1}\one=c_{\one}w^*$ and
$\one'\Sigma^{-1}E\Sigma^{-1}\one=c_{\one}^2w^{*\prime}Ew^*$ give
$\wgmv(\Sigma+E)=w^*-\Sigma^{-1}Ew^*+c_{\one}(w^{*\prime}Ew^*)w^*+O(\norm E^2)=w^*-MEw^*+O(\norm E^2)$.
By \cref{lem:quad-loss}(a) and $M\Sigma M=M$ (which follows from $\one'M=0$),
$V(\wgmv(\Sigma+E))-V^*=(Ew^*)'M(Ew^*)+O(\norm E^3)$. Next,
$M=\Sigma^{-1/2}(I-uu'/u'u)\Sigma^{-1/2}$ is positive semidefinite with kernel
$\Sigma^{1/2}\operatorname{span}(u)=\operatorname{span}(\one)$. Finally $Ew^*=c\one$ iff
$(\Sigma+E)w^*=(1/c_{\one}+c)\one$, and since $w^{*\prime}(\Sigma+E)w^*=1/c_{\one}+c>0$, this holds iff
$w^*=\wgmv(\Sigma+E)$.
\end{proof}

\begin{proof}[Proof of \cref{cor:gmv-drivers}(i)]
For a stationary ergodic generator with finite second moments, the ergodic theorem applied to
$r_t$ and $r_tr_t'$ gives $\hat\Sigma_{\mathrm{syn}}\to\Sigma_G$ almost surely as $T_s\to\infty$.
$\hat\Sigma\mapsto V(\wgmv(\hat\Sigma))$ is continuous on positive definite matrices, so
$V(\wgmv(\hat\Sigma_{\mathrm{syn}}))/V^*\to V(\wgmv(\Sigma_G))/V^*$, which is bounded by
\cref{prop:gmv-sensitivity}(b) with $\kappa_G$ and depends on $\Sigma_G$ only up to scale by part
(a). Any other feature of the generator enters the limit only through $\Sigma_G$.
\end{proof}

Parts (ii) and (iii) of \cref{cor:gmv-drivers} are heuristic: they substitute
$\hat\Sigma_{\mathrm{syn}}=\Sigma_G+O_\Prob(T_s^{-1/2})$ into \cref{prop:gmv-sensitivity}(d) and
keep the leading bias and variance terms, neglecting higher-order terms and mean estimation.

\begin{proof}[Proof of \cref{prop:no-free-lunch}]
Let $Y=g(R_{\mathrm{tr}},U)$ with $U$ independent of $(R_{\mathrm{tr}},\Theta)$, where $\Theta$
is any \TEST{} quantity, and let $\pi(Y)$ be integrable with values in a convex feasible set
$\mathcal W\subset\R^N$. By independence,
$\E[\pi(Y)\mid R_{\mathrm{tr}},\Theta]=\int\pi(g(R_{\mathrm{tr}},u))\,\Prob_U(du)=\E[\pi(Y)\mid R_{\mathrm{tr}}]=\bar w$,
a function of $R_{\mathrm{tr}}$ only. The mean of a $\mathcal W$-valued integrable vector lies in
the convex set $\mathcal W$, so $\bar w$ is feasible. For a loss $\ell(w;\Theta)$ convex in $w$,
the conditional Jensen inequality given $(R_{\mathrm{tr}},\Theta)$ yields
$\ell(\bar w;\Theta)\le\E[\ell(\pi(Y);\Theta)\mid R_{\mathrm{tr}},\Theta]$.
$V(w)$ and $-\CE(w)$ are convex in $w$.
\end{proof}

\begin{proof}[Proof of \cref{ex:scale-blind}]
Windows satisfy $Y\eqd\sqrt c\,X$. The coupling $(X,\sqrt cX)$ gives
$\W(P,Q)\le\abs{\sqrt c-1}\E\norm X$, and the 1-Lipschitz function $\pm\norm\cdot$ attains this
value, so $\W(P,Q)=\abs{\sqrt c-1}\E\norm X>0$ for $c\neq1$. The standardisation is a bijection, so
the standardised laws differ too, and under \cref{ass:consistency} the critic test rejects with
probability tending to one (\cref{prop:consistency}). The GMV loss of $\Sigma_G=c\Sigma$ is zero by
\cref{prop:gmv-sensitivity}(a).
\end{proof}

\begin{proof}[Proof of \cref{ex:spread}]
With $\Sigma=v_1v_1'+\eta^2v_2v_2'$ and $\Sigma_G=v_1v_1'+s\eta^2v_2v_2'$ ($0<s<1$),
$\Sigma^{-1/2}=v_1v_1'+\eta^{-1}v_2v_2'$ and $A_\Sigma=v_1v_1'+s\,v_2v_2'$, so $\kappa=1/s$. With
$u=\Sigma^{-1/2}\one=(v_1'\one)v_1+\eta^{-1}(v_2'\one)v_2$,
$A_\Sigma^{-1/2}u=(v_1'\one)v_1+s^{-1/2}\eta^{-1}(v_2'\one)v_2$ has equal mass on $v_1,v_2$ iff
$\abs{v_2'\one}/\abs{v_1'\one}=\eta\sqrt s$. Since $v_1'\one=\cos\phi+\sin\phi$ and
$v_2'\one=\cos\phi-\sin\phi$, for $\phi\in(0,\pi/4)$ this reads
$(1-\tan\phi)/(1+\tan\phi)=\eta\sqrt s$, equivalently $\tan\phi=(1-\eta\sqrt s)/(1+\eta\sqrt s)$,
as stated (the right-hand side lies in $(0,1)$ because $0<\eta\sqrt s<1$). By
\cref{lem:kantorovich-ineq} the ratio equals $(1+\kappa)^2/(4\kappa)=(1+s)^2/(4s)$, which is
$1.5625$ at $s=1/4$ for every $\eta$. For \iid{} Gaussian windows of length $L$, the coupling
$Y_t=\Sigma_G^{1/2}\Sigma^{-1/2}X_t$, day by day, gives
$\E\norm{X-Y}^2=L\Tr\{(\Sigma^{1/2}-\Sigma_G^{1/2})^2\}=L\eta^2(1-\sqrt s)^2$, because the two
matrices share their eigenvectors. Hence $\W\le\{\E\norm{X-Y}^2\}^{1/2}=\sqrt L\,\eta(1-\sqrt s)\to0$
as $\eta\to0$, whatever $\phi$. The per-asset standardisation of the critic multiplies this
bound by at most $\norm{\Lambda_*^{-1}}_{\mathrm{op}}$; the two diagonal entries of the pooled
covariance $\tfrac12(\Sigma+\Sigma_G)$ tend to $1/2$ as $\eta\to0$ (because $\phi\to\pi/4$), so
$\norm{\Lambda_*^{-1}}_{\mathrm{op}}\to\sqrt2$; it is about $1.48$ at $\eta=0.1$, $s=1/4$, where
$w^*=(-1.48,2.48)$ and $\wgmv(\Sigma_G)=(-4.49,5.49)$.
\end{proof}

\begin{proof}[Proof of \cref{ex:corr-blind}]
If $G$ reproduces each asset's univariate window law exactly and simulates the assets
independently, every per-ticker null holds, so every valid univariate test has level at most
$\alpha$. Its covariance is $\Sigma_G=\diag(\Sigma)$. For two assets with standard deviations
$\sigma_1\neq\sigma_2$ and correlation $\varrho$, $\wgmv(\Sigma_G)\propto(\sigma_1^{-2},\sigma_2^{-2})$
does not depend on $\varrho$, so $V(\wgmv(\Sigma_G))$ stays bounded away from zero as
$\varrho\to1$, whereas
$V^*=\sigma_1^2\sigma_2^2(1-\varrho^2)/(\sigma_1^2+\sigma_2^2-2\varrho\sigma_1\sigma_2)\to0$. The
ratio $V(\wgmv(\Sigma_G))/V^*$ is therefore unbounded.
\end{proof}

\begin{proof}[Proof of \cref{ex:drift}]
Consider the daily returns of one asset, \iid{} Gaussian with known $\sigma$, mean $\mu$ under the
real law and $\mu+\delta_\mu$ under the synthetic law; a window test is a function of the daily
returns it covers. By \cref{prop:oracle-bound}(b), with the null point at which both samples have
the synthetic law, a two-sample test of level $\alpha$ that uses $\Tte$ real days has power at
most $\bNP$ for $N(\mu+\delta_\mu,\sigma^2)^{\otimes\Tte}$ against $N(\mu,\sigma^2)^{\otimes\Tte}$,
which is $1-\Phi(z_{1-\alpha}-\abs{\delta_\mu}\sqrt{\Tte}/\sigma)$, however many synthetic draws it
uses. With $\abs{\delta_\mu}=\sigma/\sqrt{\Ttr}$, the standard error of the \TRAIN{} mean, and
$\Tte/\Ttr=615/1258$, the noncentrality is $\sqrt{615/1258}\approx0.70$ and the bound is
$1-\Phi(1.645-0.70)\approx0.17$. (If the other assets are observed, correlated with this one and
correctly specified, the noncentrality becomes $\abs{\delta_\mu}\{\Tte(\Sigma^{-1})_{ii}\}^{1/2}$,
which is larger; the bound concerns the univariate comparison.) For the utility, the
budget-constrained mean-variance weights are
$\wmv(\hat\mu,\Sigma)=w^*+\gamma^{-1}M\hat\mu$, as the first-order condition of
\cref{lem:quad-loss}(c) shows. Hence $\wmv(\hat\mu,\Sigma)-\wmv(\mu,\Sigma)=\gamma^{-1}M(\hat\mu-\mu)$,
and by \cref{lem:quad-loss}(c) and $M\Sigma M=M$ the certainty-equivalent loss is
$(2\gamma)^{-1}(\hat\mu-\mu)'M(\hat\mu-\mu)$.
\end{proof}
